\documentclass[11pt]{amsart}
\usepackage{amsmath,amssymb,amsthm,mathtools}
\usepackage{geometry}
\usepackage{booktabs}
\usepackage{xcolor}
\usepackage{hyperref}
\usepackage{enumitem}
\geometry{margin=1in}
\hypersetup{colorlinks=true,linkcolor=blue!70!black,citecolor=blue!70!black,urlcolor=blue!70!black,
  pdftitle={A Lyapunov theorem for the de Bruijn-Newman flow},
  pdfauthor={David Alejandro Trejo Pizzo}}

\newtheorem{theorem}{Theorem}[section]
\newtheorem{proposition}[theorem]{Proposition}
\newtheorem{lemma}[theorem]{Lemma}
\newtheorem{corollary}[theorem]{Corollary}
\newtheorem{conjecture}[theorem]{Conjecture}
\theoremstyle{definition}
\newtheorem{definition}[theorem]{Definition}
\theoremstyle{remark}
\newtheorem{remark}[theorem]{Remark}
\newtheorem{observation}[theorem]{Observation}

\newcommand{\re}{\operatorname{Re}}
\newcommand{\im}{\operatorname{Im}}
\newcommand{\half}{\tfrac12}
\newcommand{\Lyap}{\mathcal L}
\newcommand{\Norm}[1]{\lVert #1\rVert}

\title[A de Bruijn--Newman Lyapunov Theorem]{A Lyapunov Theorem for the de Bruijn--Newman Flow,\\[3pt]
a Dynamical No-Go, and the Extremal Form of the\\[2pt]
Riemann Hypothesis as Gap Universality}
\author{David Alejandro Trejo Pizzo}
\thanks{Independent Researcher. \texttt{dtrejopizzo@gmail.com}}


\begin{document}
\nocite{bgNewman,bgCSV}
\begin{abstract}
We study the Riemann Hypothesis through the de Bruijn--Newman heat flow, for which $\mathrm{RH}\iff\Lambda=0$
(using Rodgers--Tao's $\Lambda\ge0$ and Polymath15's $\Lambda\le0.2$). We record three RH-independent results.
(i) A clean \emph{Lyapunov theorem}: along the Coulomb-gas zero dynamics, the total off-line-ness
$\Lyap=\sum_j(\im z_j)^2$ obeys $\dot\Lyap=-2\sum_{j\ne k}(y_j-y_k)^2/|z_j-z_k|^2\le 0$, so $\Lyap$ is a strict
Lyapunov function and the critical line is the flow's global attractor; as a by-product any hypothetical
off-line zero sits at imaginary height $\le\sqrt{2\Lambda}\le0.63$. (ii) A \emph{dynamical no-go} (N5): the
identity for $\dot\Lyap$ is \emph{arithmetic-blind} (it holds for every conjugate-symmetric configuration), so
no argument using the flow alone can prove RH; in particular this subsumes the entire class of
``critical-line-as-attractor'' dissipative-dynamics / PT-symmetry heuristics. (iii) The \emph{extremal
reduction}: $\Lambda$ equals the first backward-collision time of the zeros under the flow, an extremal (not
average) quantity, and a single off-line zero already forces $\Lambda>0$. Measuring real $\zeta$ zeros we find
the normalized backward-collision margin $t_c/s^2$ clusters at a universal negative constant $\approx-0.02$
(the sine-kernel prediction), giving the sharpest statement of the target: \emph{RH follows from the
unconditional GUE gap universality of the zeros}. We locate, precisely and quantitatively, the single
irreducible arithmetic statement at which the dynamical approach---and, as a companion map shows, every
current analytic approach---arrives.
\end{abstract}

\maketitle



\tableofcontents
\newpage

\section{Introduction}\label{sec:intro}
The de Bruijn--Newman constant $\Lambda$ is defined through the one-parameter deformation
\begin{equation}\label{eq:Ht}
H_t(z)=\int_0^\infty e^{tu^2}\,\Phi(u)\,\cos(zu)\,du,\qquad H_0(z)\propto\xi(\half+\tfrac{iz}{2}),
\end{equation}
where $\xi$ is the Riemann $\xi$-function and $\Phi$ its (super-exponentially decaying) kernel; $H_t$ satisfies
the backward heat equation $\partial_t H_t=-\partial_{zz}H_t$. By the theorems of de Bruijn and Newman, there
is a constant $\Lambda\in\mathbb R$ such that $H_t$ has only real zeros iff $t\ge\Lambda$. Newman conjectured
$\Lambda\ge0$ (``RH, if true, is only barely so''); this was proved by Rodgers and Tao \cite{RT}, while
Polymath15 \cite{P15} established $\Lambda\le0.2$. Consequently
\begin{equation}\label{eq:RHL}
\mathrm{RH}\ \Longleftrightarrow\ \Lambda\le 0\ \Longleftrightarrow\ \Lambda=0.
\end{equation}
This reframes RH as the statement that the flow's zeros are \emph{already} all real at $t=0$, the boundary of
the all-real regime. It is attractive because the reality of the zeros for $t\ge\Lambda>0$ is a \emph{theorem},
not RH, and the only remaining content is the propagation of reality down to $t=0$, governed by the zero
dynamics rather than by the explicit formula.

This paper records what that dynamical viewpoint yields rigorously. The results are RH-independent: a Lyapunov
theorem and its attractor corollary (\S\ref{sec:lyap}), a no-go showing the flow alone cannot decide RH
(\S\ref{sec:nogo}), and the extremal reduction of $\Lambda$ to a backward-collision time together with the
numerical observation that pins the target as gap universality (\S\ref{sec:extremal}). The companion paper
\cite{P11} places these alongside the static reformulations; here we develop the dynamical side in full.

\paragraph{The zero dynamics.} Where $H_t$ has simple zeros $z_j(t)$, differentiating $H_t(z_j(t))=0$ and using
$\partial_tH=-\partial_{zz}H$ gives the Calogero/Coulomb flow
\begin{equation}\label{eq:flow}
\dot z_j\;=\;2\sum_{k\ne j}\frac{1}{z_j-z_k},
\end{equation}
the classical zero dynamics of the heat flow \cite{CSV}. (We normalize the time so the constant is $2$; all
statements below are invariant under rescaling $t$.) Since $H_t$ has real coefficients, the zeros are symmetric
under conjugation; complex zeros occur in conjugate pairs.

\section{A Lyapunov theorem and the attractor}\label{sec:lyap}

Write $z_j=x_j+iy_j$ and define the \emph{total off-line-ness}
\begin{equation}\label{eq:L}
\Lyap\;:=\;\sum_j (\im z_j)^2\;=\;\sum_j y_j^2\;\ge0,\qquad \Lyap=0\iff\text{all zeros real.}
\end{equation}

\begin{theorem}[Lyapunov monotonicity]\label{thm:lyap}
Along the flow \eqref{eq:flow},
\begin{equation}\label{eq:Ldot}
\dot\Lyap\;=\;-2\sum_{j\ne k}\frac{(y_j-y_k)^2}{|z_j-z_k|^2}\;\le\;0,
\end{equation}
with equality iff all $y_j$ are equal; for a conjugate-symmetric configuration this forces all $y_j=0$. Hence
$\Lyap$ is a strict Lyapunov function for the forward flow, and the set of all-real configurations is its
global attractor.
\end{theorem}

\begin{proof}
$\dot\Lyap=2\sum_j y_j\dot y_j=2\sum_j y_j\,\im\dot z_j$. By \eqref{eq:flow},
$\im\dot z_j=2\sum_{k\ne j}\im(z_j-z_k)^{-1}=-2\sum_{k\ne j}(y_j-y_k)/|z_j-z_k|^2$, so
$\dot\Lyap=-4\sum_{j\ne k}y_j(y_j-y_k)/|z_j-z_k|^2$. Symmetrizing under $j\leftrightarrow k$ (each unordered pair
twice) replaces $y_j(y_j-y_k)$ by $\half[(y_j-y_k)^2]$, giving \eqref{eq:Ldot}. Each summand is $\ge0$;
equality needs $y_j=y_k$ for all $j,k$, and conjugate symmetry ($\sum y_j=0$, $y$'s in $\pm$ pairs) then forces
$y_j\equiv0$.
\end{proof}

This is the rigorous content of the intuition that ``the critical line is a stable attractor'': under the heat
flow, off-line zeros are pulled toward the real axis, monotonically in the quadratic measure $\Lyap$. We record
two consequences.

\begin{corollary}[Healing rate; height bound]\label{cor:height}
An isolated conjugate pair at $\pm iy_0$ obeys $\frac{d}{dt}y_0^2=-2$ exactly, hence heals (reaches the real
axis) by time $y_0^2/2$. Therefore any zero of $H_0=\xi$ off the real axis (i.e.\ any RH-violating zero) has
imaginary part $|y_0|\le\sqrt{2\Lambda}\le\sqrt{0.4}\approx0.63$.
\end{corollary}

\begin{proof}
For the isolated pair, \eqref{eq:Ldot} gives $\dot\Lyap=-2\cdot 2\cdot\frac{(2y_0)^2}{(2y_0)^2}=-4$ and
$\Lyap=2y_0^2$, so $\frac{d}{dt}y_0^2=-2$ and $y_0^2(t)=y_0^2(0)-2t$. A pair present at $t=0$ heals by
$t=y_0^2(0)/2$; since the all-real regime begins at $t=\Lambda$, $y_0^2(0)/2\le\Lambda$, and $\Lambda\le0.2$.
\end{proof}

\section{A dynamical no-go: the flow is arithmetic-blind}\label{sec:nogo}

Theorem~\ref{thm:lyap} is appealing, but it cannot prove RH, for two structural reasons that we make precise.

First, \emph{wrong direction}: forward in $t$, $\Lyap$ decreases to $0$ (the regime $t\ge\Lambda$). RH concerns
$t=0\le\Lambda$, reached by flowing \emph{backward}, where $\Lyap$ \emph{increases}; the monotonicity then
yields only $\Lyap(0)\ge\Lyap(\Lambda)=0$, which is trivially true and not $\Lyap(0)=0$.

Second, and decisively, \emph{arithmetic-blindness}.

\begin{theorem}[Dynamical no-go, N5]\label{thm:nogo}
The identity \eqref{eq:Ldot} holds for \emph{every} finite conjugate-symmetric configuration $\{z_j\}$,
independently of any arithmetic origin. Consequently no argument that uses only the flow \eqref{eq:flow}
together with a Lyapunov/monotonicity functional of the zero positions can distinguish $\xi$ from an arbitrary
such function, and hence cannot prove RH: it would prove ``every function in the de Bruijn--Newman class
satisfies RH,'' which is false (generic initial data have $\Lambda>0$). Any RH attack based on a dissipative
dynamics with the critical line as attractor is of this form, and is subject to this obstruction.
\end{theorem}

\begin{proof}
The derivation of \eqref{eq:Ldot} used only \eqref{eq:flow} and $z_j=x_j+iy_j$, never the values $\{z_j(0)\}$.
The flow \eqref{eq:flow} is the same vector field for all initial data. A functional $F(\{z_j\})$ of the
positions therefore has $\dot F$ determined by the universal field; only the initial value $F(\{z_j(0)\})$
carries the arithmetic, and that value is exactly what RH asserts. There exist functions in the class with
$\Lambda>0$ (off-line zeros), e.g.\ suitable Fourier transforms violating the Hermite--Biehler condition; for
these the same $\dot\Lyap\le0$ holds. Hence the flow cannot separate $\Lambda\le0$ from $\Lambda>0$.
\end{proof}

\begin{remark}
Theorem~\ref{thm:nogo} explains, in a single sentence, why the heat flow does not escape the difficulty even
though it introduces a genuinely new parameter $t$: the new degree of freedom is arithmetic-blind. RH lives
entirely in the \emph{initial condition} $H_0=\xi$, i.e.\ in where the zeros start and how they are spaced. To
use the flow one must re-inject $\zeta$'s zero statistics; \S\ref{sec:extremal} identifies exactly which
statistic.
\end{remark}

\section{The extremal reduction, and the target as gap universality}\label{sec:extremal}

Because the flow preserves the real axis forward but contracts gaps backward, $\Lambda$ admits a clean
extremal description.

\begin{proposition}[$\Lambda$ as first backward collision]\label{prop:collision}
Under \eqref{eq:flow} a pair of real zeros with gap $g$ obeys $g(t)^2=g_0^2+8t$ to leading (two-body) order, so
backward in $t$ the gap shrinks and the pair collides (forming a double zero, then a conjugate pair) at
$t_c=-g_0^2/8$. The de Bruijn--Newman constant is the supremum of such collision times over all pairs and
heights: $\Lambda=\sup t_c$. In particular a \emph{single} off-line zero forces $\Lambda>0$; $\Lambda$ is an
\emph{extremal}, not an averaged, functional of the zeros.
\end{proposition}

This is the crux. Any control of RH through the flow must bound the \emph{worst} pair, uniformly in height ---
no averaged statistic (mean spacing, integrated pair correlation) suffices, since each is insensitive to a
sparse set of anomalously tight pairs, while one such pair decides $\Lambda$.

\paragraph{A universal collision margin (numerical).} We backward-flow blocks of consecutive real $\zeta$ zeros
(via \texttt{mpmath.zetazero}) until the tightest pair collides, recording the collision time $t_c$ and the
local mean spacing $s\sim 2\pi/\log T$:

\begin{center}
\begin{tabular}{rrrrr}
\toprule
height $T$ & min gap & mean spacing $s$ & $t_c$ & $t_c/s^2$\\
\midrule
$99$ & 0.845 & 2.281 & $-0.105$ & $-0.0202$\\
$850$ & 0.484 & 1.280 & $-0.032$ & $-0.0195$\\
$2547$ & 0.425 & 1.046 & $-0.025$ & $-0.0229$\\
\bottomrule
\end{tabular}
\end{center}

In absolute units the margin $|t_c|$ shrinks with height (the spacing $s\to0$), but in mean-spacing units
$t_c/s^2$ clusters at a \emph{universal negative constant} $\approx-0.02$. This is the sine-kernel / GUE
prediction for the (normalized) closest-pair collision time, and it is consistent with the marginal value
$\Lambda=0$: the supremum of $t_c$ tends to $0^-$ in absolute units while the normalized margin stays strictly
below $0$. (The computation is a deliberate probe---three blocks, a windowed flow without the mean field---so
the constant is suggestive, not established.)

\begin{observation}[The sharpest target]\label{obs:target}
By Proposition~\ref{prop:collision}, RH $=\{\Lambda\le0\}$ follows if the normalized backward-collision margin
$t_c/s^2$ stays a negative constant uniformly in $T$ --- equivalently, if the closest-pair gap statistics of
the $\zeta$ zeros are GUE-universal \emph{unconditionally}. This is the extremal, quantitative form of the
target. Its uniform control is precisely the open, currently RH-conditional, statement of pair-correlation /
gap universality (Montgomery's conjecture is proved only in a restricted range \cite{Mont}).
\end{observation}

\section{Conclusion}\label{sec:concl}

The de Bruijn--Newman flow gives a clean dynamical picture of RH and three RH-independent results: a Lyapunov
theorem with the critical line as attractor (Theorem~\ref{thm:lyap}) and the height bound
(Corollary~\ref{cor:height}); a dynamical no-go showing the flow alone is arithmetic-blind and subsuming the
PT-symmetry/dissipative heuristics (Theorem~\ref{thm:nogo}); and the extremal reduction of $\Lambda$ to a
backward-collision time, with the numerical universality of the normalized margin pinning the target as the
\emph{unconditional gap universality} of the zeros (Proposition~\ref{prop:collision},
Observation~\ref{obs:target}). The viewpoint does not cross the wall; it compresses RH, with unusual precision,
into a single extremal arithmetic statement. Establishing that statement unconditionally---the GUE universality
of the closest-pair gaps, uniform in height---is the new analytic input the problem demands, and the subject of
a separate program.

\begin{thebibliography}{9}
\bibitem{RT} B.~Rodgers and T.~Tao, \emph{The de Bruijn--Newman constant is non-negative}, Forum Math. Pi \textbf{8} (2020), e6.
\bibitem{P15} D.~H.~J.~Polymath, \emph{Effective approximation of heat flow evolution of the Riemann $\xi$ function, and a new upper bound for the de Bruijn--Newman constant}, Res. Math. Sci. \textbf{6} (2019), 31.
\bibitem{CSV} G.~Csordas, W.~Smith, R.~S.~Varga, \emph{Lehmer pairs of zeros, the de Bruijn--Newman constant $\Lambda$, and the Riemann hypothesis}, Constr. Approx. \textbf{10} (1994), 107--129.
\bibitem{Mont} H.~L.~Montgomery, \emph{The pair correlation of zeros of the zeta function}, Proc. Sympos. Pure Math. \textbf{24} (1973), 181--193.
\bibitem{P11} D.~A.~Trejo Pizzo, \emph{The band-limited Weil--Carleson form and a five-language map of the Riemann wall}, preprint, 2026.
\bibitem{bgNewman} C.~M.~Newman, \emph{Fourier transforms with only real zeros}, Proc.\ Amer.\ Math.\ Soc.\ \textbf{61} (1976), 245--251.
\bibitem{bgCSV} G.~Csordas, W.~Smith, and R.~S.~Varga, \emph{Lehmer pairs of zeros, the de Bruijn--Newman constant $\Lambda$, and the Riemann Hypothesis}, Constr.\ Approx.\ \textbf{10} (1994), 107--129.
\end{thebibliography}

\end{document}
